\documentclass[11pt,reqno]{amsart}
\usepackage[T1]{fontenc}
\usepackage{lmodern}
\usepackage[letterpaper,margin=1in]{geometry}
\usepackage{microtype}
\usepackage{amsmath,amssymb,mathtools}
\usepackage[hidelinks,pdfusetitle]{hyperref}
\usepackage{enumitem}
\usepackage{needspace}
\numberwithin{equation}{section}
\newtheorem{theorem}{Theorem}[section]
\newtheorem{proposition}[theorem]{Proposition}
\newtheorem{lemma}[theorem]{Lemma}
\newtheorem{corollary}[theorem]{Corollary}
\theoremstyle{definition}
\newtheorem{definition}[theorem]{Definition}
\theoremstyle{remark}
\newtheorem{remark}[theorem]{Remark}
\newcommand{\R}{\mathbb R}
\newcommand{\Z}{\mathbb Z}
\newcommand{\N}{\mathbb N}
\newcommand{\E}{\mathbb E}
\newcommand{\1}{\mathbf 1}
\newcommand{\Conf}{\operatorname{Conf}}
\newcommand{\Var}{\operatorname{Var}}
\newcommand{\PV}{\operatorname{PV}}
\newcommand{\Ent}{\operatorname{Ent}}
\newcommand{\Sine}{\operatorname{Sine}}
\newcommand{\dd}{\,d}
\newcommand{\gap}{\mathfrak g}
\newcommand{\calB}{\mathcal B}
\newcommand{\calI}{\mathcal I}
\newcommand{\calE}{\mathcal E}
\newcommand{\calF}{\mathcal F}
\newcommand{\calG}{\mathcal G}
\newcommand{\norm}[1]{\left\lVert #1\right\rVert}
\newcommand{\abs}[1]{\left\lvert #1\right\rvert}
\DeclareMathOperator{\supp}{supp}
\DeclareMathOperator{\diver}{div}
\allowdisplaybreaks[2]
\setlength{\emergencystretch}{1.5em}
\setlist[enumerate]{leftmargin=*,itemsep=3pt,topsep=5pt}
\title[Uniqueness of the logarithmic DLR state]{Uniqueness of the stationary logarithmic DLR state}
\date{Research manuscript draft\\9 September 2026}
\subjclass[2020]{60G55, 60K35, 82B21, 60B20}
\keywords{Logarithmic gas, DLR equations, Sine beta process, hyperuniformity, Palm measure, entropy}

\begin{document}
\begin{abstract}
We give an entropy comparison argument for uniqueness of the stationary
one-dimensional logarithmic DLR state at every inverse temperature
$\beta>0$. The comparison theorem is first formulated under explicit count
variance assumptions: a linear variance bound, sublinear variance growth,
and integrability of the variance against $t^{-2}\dd t$. Two conditional
laws with the same number of particles are compared with one deterministic
lattice-boundary law carrying a weak quadratic confinement. A static
concentration estimate for averages of finite gap vectors bounds their
expectation difference by relative entropy per particle. Integrated
discrepancy control makes the boundary contribution to the Fisher
information tight, while sublinear variance makes the confinement cost
negligible. We then derive the required variance assumptions from finite
logarithmic electric-field energy by stationarization and a covariance
spectral estimate. With the standard regularized interpretation of the
renormalized energy, this identifies the stationary finite-energy solution
of the canonical DLR equations with $\Sine_\beta$.
\end{abstract}
\maketitle

\noindent\textit{Manuscript status.}
This is a working draft prepared for mathematical review. The argument is
developed from the three supplied preprints listed in the bibliography.
No external literature search or claim of priority is made. The versions
actually used, including the energy convention discussed in
Section~\ref{sec:energy}, are specified explicitly.

\section{Introduction and statements}

The process $\Sine_\beta$ describes the microscopic bulk limit of
one-dimensional logarithmic gases. Dereudre, Hardy, Lebl\'e and
Ma\"ida~\cite{DHLM} show that it satisfies canonical Dobrushin--Lanford--Ruelle
equations and that every stationary finite-energy solution of those
equations is number-rigid. Their Section~1.3 asks whether the stationary
solution is unique. The conditional universality theorem of Kuijlaars and
Mi\~na-D\'iaz~\cite{KMD} supplies a deterministic route at $\beta=2$.
For general $\beta$, Bourgade, Erd\H{o}s and Yau~\cite{BEY} develop a
comparison method for local logarithmic Gibbs measures, together with
strong particle-location estimates for convex analytic potentials.

The argument below uses the convexity and entropy ingredients of that
method. It compares two candidate infinite-volume laws with the same
finite reference law. It does not require a limiting theorem for the
reference law, nor uniform microscopic rigidity of every frozen boundary.
The finite-energy input is instead used to obtain an integrated
discrepancy estimate.

We write $\gamma$ for a locally finite simple counting measure on $\R$,
$N_I(\gamma)=\gamma(I)$, and $\theta_a\gamma=\gamma-a$. All stationary
processes in the comparison theorem have intensity one. For such a
process $P$, put
\[
v_P(t)=\Var_P N_{(0,t]},\qquad t>0.
\]
The canonical DLR property means that, conditionally on the exterior of a
bounded interval and on its particle number, the interior points have the
normalized logarithmic Gibbs density at inverse temperature $\beta$.
The infinite exterior interaction is defined by symmetric cutoffs and
relative energies, as in \cite[Section~2.1]{DHLM}.

\Needspace{10\baselineskip}
\begin{theorem}[Uniqueness under discrepancy control]\label{thm:comparison}
Fix $\beta>0$. There is at most one stationary simple point process of
intensity one satisfying the canonical logarithmic DLR equations and the
following three conditions:
\begin{align}
v_P(t)&\le C_Pt \quad(t>0),\label{eq:linear}\\
v_P(t)&=o(t)\quad(t\to\infty),\label{eq:hyper}\\
\int_1^\infty v_P(t)t^{-2}\dd t&<\infty.\label{eq:integrated}
\end{align}
The constant $C_P$ need not be uniform over candidate processes.
\end{theorem}

For clarity, all three conditions are retained in the statement.
Section~\ref{sec:energy} proves them together for each ergodic finite-energy
DLR solution. The comparison proof itself requires no ergodicity.

\begin{theorem}[Finite-energy uniqueness]\label{thm:finite-energy}
Fix $\beta>0$ and interpret the renormalized logarithmic energy with the
standard inner electric-field truncation specified in
Section~\ref{sec:energy}. If $P$ is stationary, satisfies the canonical DLR
equations of \cite[Definition~3.1]{DHLM}, and $\E_P W<\infty$, then
\[
P=\Sine_\beta.
\]
\end{theorem}

Theorem~\ref{thm:finite-energy} follows from
Theorem~\ref{thm:comparison}, the electric-field argument, and the ergodic
DLR decomposition in \cite[Proposition~3.13]{DHLM}. It makes no assertion
about unnormalized conditional-density specifications allowing different
intensities.

\subsection{Organization and main estimates}

Section~\ref{sec:entropy} proves a static comparison inequality for
averages of consecutive gap vectors:
\[
|\mu(F_K)-\omega_{K,T}(F_K)|^2
\le \frac{C_GT}{K}I(\mu\mid\omega_{K,T}).
\]
Here $\omega_{K,T}$ is a lattice-boundary logarithmic gas with quadratic
confinement of strength $T^{-1}$. Section~\ref{sec:palm} collects the
counting and Palm facts needed for random particle-ended intervals.
Sections~\ref{sec:blocks} and \ref{sec:forces} control, respectively, the
confinement cost and the exterior-force contribution. These estimates
are assembled in Section~\ref{sec:proof-main}. The electric-field
argument in Section~\ref{sec:energy} verifies the hypotheses of
Theorem~\ref{thm:comparison} for finite-energy states.

\section{A static comparison inequality for gap statistics}
\label{sec:entropy}

Fix $K\ge1$, write $L=K+1$, and let
\[
\Omega_K=\{u\in\R^K:0<u_1<\cdots<u_K<L\}.
\]
Set $u_0=0$, $u_{K+1}=L$, and $g_i(u)=u_{i+1}-u_i$ for
$0\le i\le K$. The reference exterior is
\[
\xi_K^{\rm lat}=\{j\in\Z:j\le0\text{ or }j\ge L\}.
\]
With any fixed $r_*\in(0,L)$, define
\begin{equation}\label{eq:lattice-energy}
H_K^{\rm lat}(u)=
-\sum_{1\le i<j\le K}\log(u_j-u_i)
-\sum_{i=1}^K\PV\sum_{p\in\xi_K^{\rm lat}}
       \log\abs{\frac{u_i-p}{r_*-p}}.
\end{equation}
The reference changes only an additive constant. Equivalently, Euler's
product for the gamma function gives the completely finite expression
\[
H_K^{\rm lat}(u)=
-\sum_{i<j}\log(u_j-u_i)
+\sum_{i=1}^K\{\log\Gamma(u_i)+\log\Gamma(L-u_i)\}
+\text{constant}.
\]
The paired exterior series converges, and its derivatives of order at
least two converge locally uniformly. The endpoint singularities are
$-\log u_i-\log(L-u_i)$.
For $T>0$ put
\begin{equation}\label{eq:reference}
\omega_{K,T}(\dd u)=Z_{K,T}^{-1}e^{-\Phi_{K,T}(u)}\dd u,\qquad
\Phi_{K,T}(u)=\beta H_K^{\rm lat}(u)
                  +\frac1{2T}\sum_{i=1}^K(u_i-i)^2.
\end{equation}
The normalizing constant is finite and strictly positive.

Fix an integer $r\ge1$ and a function $G\in C_c^2((0,\infty)^r)$.
For $K>r$, define
\begin{equation}\label{eq:gap-observable}
F_K(u)=\frac1{K-r}
\sum_{i=1}^{K-r}G(g_i(u),\ldots,g_{i+r-1}(u)).
\end{equation}
This uses only gaps between interior particles. Omitting finitely many
edge windows has no effect on the limiting average.

For probability measures $\mu\ll\omega$ on $\Omega_K$, our entropy and
Fisher information conventions are
\[
S(\mu\mid\omega)=\int\log\frac{\dd\mu}{\dd\omega}\dd\mu,\qquad
I(\mu\mid\omega)=
\int\abs{\nabla\log\frac{\dd\mu}{\dd\omega}}^2\dd\mu.
\]
Fisher information is interpreted in the weak Sobolev sense, with value
$+\infty$ outside its domain.

\begin{proposition}[Gap entropy bound]\label{prop:gap-entropy}
For fixed $\beta,r,G$, there is a constant $C_G$ independent of $K,T$
such that, for $K\ge2r+1$,
\begin{align}
|\mu(F_K)-\omega_{K,T}(F_K)|^2
&\le \frac{C_G}{K}S(\mu\mid\omega_{K,T}),\label{eq:entropy-gap}\\
|\mu(F_K)-\omega_{K,T}(F_K)|^2
&\le \frac{C_GT}{K}I(\mu\mid\omega_{K,T}).\label{eq:fisher-gap}
\end{align}
\end{proposition}

\begin{proof}
Write $\Phi=\Phi_{K,T}$. For $v\in\R^K$, extend $v$ by
$v_0=v_{K+1}=0$. Nearest-neighbor interactions and the two endpoint
charges give
\begin{equation}\label{eq:hessian}
v^\top\Phi''(u)v
\ge\beta\sum_{i=0}^K
\frac{(v_{i+1}-v_i)^2}{g_i(u)^2}+\frac1T|v|^2.
\end{equation}
The remaining logarithmic interactions have nonnegative Hessians.

Regard $F_K$ as a function of the gaps and set
$q_i=\partial F_K/\partial g_i$. Bounded overlap of the windows in
\eqref{eq:gap-observable} gives $|q_i|\le C_G/K$. Moreover $q_i=0$
unless $g_i$ lies in a fixed bounded interval determined by the support
of the derivatives of $G$. Consequently,
\[
\begin{split}
|\nabla F_K\cdot v|^2
&=\abs{\sum_iq_i(v_{i+1}-v_i)}^2\\
&\le
\left(\beta^{-1}\sum_iq_i^2g_i^2\right)
\left(\beta\sum_i\frac{(v_{i+1}-v_i)^2}{g_i^2}\right)
\le \frac{C_G}{K}v^\top\Phi''v.
\end{split}
\]
Thus
\begin{equation}\label{eq:gradient-metric}
\nabla F_K^\top(\Phi'')^{-1}\nabla F_K\le C_G/K.
\end{equation}
The same support property for second derivatives, followed by
$2|ab|\le a^2+b^2$ and bounded overlap, gives
\begin{equation}\label{eq:tilt-hessian}
|v^\top F_K''v|
\le \frac{C_G}{K}
v^\top(\beta H_K^{{\rm lat}\,\prime\prime})v.
\end{equation}

Let $\omega_s$ have density proportional to $e^{sF_K}\dd\omega_{K,T}$.
For $|s|\le c_GK$, \eqref{eq:tilt-hessian} yields
\[
(\Phi-sF_K)''
\ge \tfrac12\beta H_K^{{\rm lat}\,\prime\prime}+T^{-1}I
\ge\tfrac12\Phi''.
\]
The Brascamp--Lieb variance inequality and
\eqref{eq:gradient-metric} imply
$\Var_{\omega_s}F_K\le C_G/K$. Integrating the second derivative of
the logarithmic moment-generating function therefore gives
\begin{equation}\label{eq:mgf}
\log\E_{\omega_{K,T}}
\exp\{s(F_K-\omega_{K,T}(F_K))\}
\le C_Gs^2/K,\qquad |s|\le c_GK.
\end{equation}
The entropy inequality, used with either sign, shows that for
$0<s\le c_GK$,
\[
|\mu(F_K)-\omega_{K,T}(F_K)|
\le \frac{S(\mu\mid\omega_{K,T})}{s}+\frac{C_Gs}{K}.
\]
Optimizing proves \eqref{eq:entropy-gap} when the entropy is at most a
fixed multiple of $K$. Above that threshold it follows from
$|F_K|\le\norm{G}_\infty$, after enlarging the constant.

Finally, \eqref{eq:hessian} gives the logarithmic Sobolev inequality
\begin{equation}\label{eq:lsi}
S(\mu\mid\omega_{K,T})\le \frac T2 I(\mu\mid\omega_{K,T}),
\end{equation}
and hence \eqref{eq:fisher-gap}. The static convexity inequalities here
are the ones used in \cite[Sections~3 and 5]{BEY}. They apply on a
convex simplex by approximation: first impose positive lower bounds on
all gaps, use the convex-domain inequalities, and then let those bounds
decrease to zero. The boundary terms have the favorable convex sign.
Bounded tilts pass by dominated convergence, and the logarithmic Sobolev
inequality passes first for smooth bounded functions and then by
Sobolev truncation. This proof has no restriction $\beta\ge1$.
\end{proof}

\section{Discrepancy and Palm preliminaries}\label{sec:palm}

Throughout Sections~\ref{sec:palm}--\ref{sec:proof-main}, $P$ satisfies the
hypotheses of Theorem~\ref{thm:comparison}. The symbol $P^0$ denotes its
Palm distribution, and $T_*$ denotes the shift from a point at the origin
to its successor.

\begin{lemma}[Almost-sure exterior regularity]\label{lem:as-regularity}
Almost surely the configuration has density one in both directions,
and, at every finite origin $a$,
\begin{equation}\label{eq:as-integral}
\int_1^\infty\frac{
  (N_{(a,a+t]}-t)^2+(N_{[a-t,a)}-t)^2}{t^2}\dd t<\infty.
\end{equation}
The symmetric principal value
$\lim_{R\to\infty}\sum_{p\in\gamma,\,0<|p|\le R}1/p$ exists,
with a possible point at the origin omitted.
These properties hold $P^0$-almost surely as well.
\end{lemma}

\begin{proof}
By \eqref{eq:linear}, Chebyshev's inequality along $t=n^2$ and
Borel--Cantelli give $N_{(0,t]}/t\to1$; monotonic interpolation removes
the subsequence. The backward statement is identical.
Tonelli and \eqref{eq:integrated} give \eqref{eq:as-integral} at $a=0$.
Write the oriented cumulative discrepancy as $D_0(x)$ for $x\in\R$.
The discrepancy of $(a,a+t]$ equals $D_0(a+t)-D_0(a)$. A change of
variable and comparability of $t$ and $a+t$ for large $t$ prove the
integrability at every finite $a$, simultaneously on the same event.

For the principal value, put
$D_{\rm sym}(t)=N_{(0,t]}-N_{[-t,0)}$. Then
\[
\sum_{1<|p|\le R,\ p\in\gamma}\frac1p
=\frac{D_{\rm sym}(R)}R-D_{\rm sym}(1)
  +\int_1^R\frac{D_{\rm sym}(t)}{t^2}\dd t.
\]
The first term vanishes by density one. The integral converges
absolutely by Cauchy--Schwarz and \eqref{eq:as-integral}. Add the finite
sum over $0<|p|\le1$. Since the resulting good event is invariant under all finite
translations, Campbell's formula transfers it to Palm measure.
\end{proof}

\begin{lemma}[Integrable Palm count envelopes]\label{lem:palm-envelope}
Under $P^0$ define
\[
C^+=\sup_{t\ge0}\frac{N_{(0,t]}}{1+t},\qquad
C^-=\sup_{t\ge0}\frac{N_{[-t,0)}}{1+t}.
\]
Then $\E_{P^0}C^++\E_{P^0}C^-<\infty$.
\end{lemma}

\begin{proof}
For a spatially rooted sample put
$C_0=\sup_{t\ge0}N_{(0,t]}/(1+t)$.
A dyadic bound gives
\[
C_0\le N_{(0,1]}+2+
 \sum_{j\ge0}2^{-j}
 \abs{N_{(0,2^{j+1}]}-2^{j+1}}.
\]
Minkowski's inequality and \eqref{eq:linear} show
$\E_PC_0^2<\infty$. For every $p\in\gamma\cap(0,1]$,
$C^+(\theta_p\gamma)\le2C_0(\gamma)$.
Campbell's formula at intensity one therefore gives
\[
\E_{P^0}C^+
=\E_P\sum_{p\in\gamma\cap(0,1]}C^+(\theta_p\gamma)
\le2\E_P[N_{(0,1]}C_0]<\infty.
\]
Reflect the argument for $C^-$.
\end{proof}

\subsection{The first positive point and the preceding gap}

Under $P^0$ enumerate the points by
$\cdots<p_{-1}<p_0=0<p_1<\cdots$, and write
$g_i=p_{i+1}-p_i$. Palm measure is invariant under $T_*$, and
Campbell's identity gives the following version of the inspection
paradox. If $a$ is the first positive point of a spatial sample, the law
of $\theta_a\gamma$ is
\begin{equation}\label{eq:first-positive}
\widetilde P^0(\dd\gamma)=g_{-1}(\gamma)P^0(\dd\gamma).
\end{equation}
Indeed, each gap ending at a point $p$ contributes the interval of
possible spatial origins of length $g_{-1}(\theta_p\gamma)$.
In particular $\E_{P^0}g_{-1}=1$.

The event that the entire gap containing the spatial origin has length
at most $M$ becomes $\{g_{-1}\le M\}$ in \eqref{eq:first-positive}.
On this event the density relative to $P^0$ is at most $M$.
Since $T_*$ preserves $P^0$, the same domination holds after any fixed
number of successor shifts, uniformly in that number. We will use this
for the two endpoints of a block.

\begin{lemma}[Spatial block averages identify Palm gap laws]
\label{lem:identify}
For fixed $r$ and bounded $G$, the expectation under $P$ of the average
of $G$ over $K-r$ successive gap windows following the first positive
point converges to $P^0(G(g_0,\ldots,g_{r-1}))$.
\end{lemma}

\begin{proof}
Let $\calI$ be the invariant sigma-algebra of $T_*$. Birkhoff's theorem
gives convergence of these bounded averages to
$\E_{P^0}[G\mid\calI]$ under Palm measure.
Density one and $\E_{P^0}g_{-1}<\infty$ give
\[
\E_{P^0}[g_{-1}\mid\calI]
=\lim_{n\to\infty}\frac1n\sum_{j=0}^{n-1}g_{-1}\circ T_*^j=1.
\]
Use \eqref{eq:first-positive}, dominated convergence with the integrable
weight $g_{-1}$, and conditioning on $\calI$.
\end{proof}

\section{DLR intervals with a fixed particle number}\label{sec:blocks}

For a spatial sample let $a$ be the first positive point and let $b$ be
the $(K+2)$-nd positive point. The open interval $(a,b)$ contains exactly
$K$ points, written $x_1<\cdots<x_K$. Put
\[
L=K+1,\qquad \ell=b-a,\qquad c=L/\ell,\qquad
u_i=c(x_i-a).
\]
Let $\zeta=\gamma-\sum_{i=1}^K\delta_{x_i}$ denote the retained exterior.
Both endpoints, all negative points, and all points beyond $b$ are
retained.

\begin{lemma}[Random-rank DLR disintegration]\label{lem:random-dlr}
Conditionally on $\zeta$, the $x_i$ have the normalized logarithmic
Gibbs density on $a<x_1<\cdots<x_K<b$, with exterior $\zeta$.
After the affine change of variables the conditional law, denoted
$\mu_\zeta$, is a logarithmic Gibbs law on $\Omega_K$ with exterior
$c(\zeta-a)$ and endpoint charges at $0,L$.
\end{lemma}

\begin{proof}
First restrict to $b<R$ for a fixed integer $R$ and disintegrate
canonical DLR in $(0,R)$. Conditional on its exterior, particle number,
and all ordered coordinates except the $K$ middle points, the remaining
coordinates have precisely the Gibbs density in the gap between $a$
and $b$. The number in $(0,R)$ is determined by the retained
configuration and $K$. The selection of $a$ as first positive point and
of $b$ as the specified rank is preserved by all such interior
resamplings. Thus the conditioning introduces no additional factor
depending on the middle coordinates. The events $\{b<R\}$ are
retained-exterior measurable and exhaust probability one.

Equivalently, one may restrict the $K$-th Campbell density in
\cite[Theorem~3.7]{DHLM} to the ordered tuple lying between the first
two positive points of the reduced configuration.
The affine Jacobian and the constant factors in logarithmic
interactions are absorbed by normalization. Relative exterior energies
are translation and scale covariant: a change of cutoff center changes
only distant shells, whose contribution vanishes by density one and
the relative summand bound $O(1/|p|)$.
\end{proof}

Write $\gap_*=a-p_-$, where $p_-$ is the last point before the spatial
origin. This is the whole gap containing that origin. It is measurable
from $\zeta$, and $a\le\gap_*$.

\begin{lemma}[Good endpoint lengths]\label{lem:length}
One has $\ell-L=o_P(\sqrt K)$. For any two candidate processes $P,Q$,
there is a deterministic sequence $r_K\downarrow0$ such that, for each
fixed $M$ and
\begin{equation}\label{eq:good}
E_{K,M}=\{\gap_*\le M,\ |\ell-L|\le r_K\sqrt K\},
\end{equation}
\[
\limsup_{K\to\infty}P(E_{K,M}^c)\le P(\gap_*>M),
\]
and the same assertion holds for $Q$.
The event $E_{K,M}$ is retained-exterior measurable.
\end{lemma}

\begin{proof}
For fixed $\varepsilon>0$, the event that the $(K+2)$-nd positive point
exceeds $K+2+\varepsilon\sqrt K$ is a count deviation of at least
$\varepsilon\sqrt K$ at a deterministic interval of length asymptotic
to $K$. Chebyshev and \eqref{eq:hyper} make its probability tend to
zero. The lower tail is identical. Since $a$ is a finite random variable
with a law not depending on $K$, subtracting $a$ proves the assertion
for $\ell$. Choose a common slowly decreasing $r_K$ by diagonalization.
The assertions about $E_{K,M}$ follow.
\end{proof}

Define the conditional quadratic cost
\begin{equation}\label{eq:J}
J_K(\zeta)=\E_{\mu_\zeta}\sum_{i=1}^K(u_i-i)^2.
\end{equation}

\begin{lemma}[Negligible quadratic cost]\label{lem:quadratic}
For every fixed $M$,
\begin{equation}\label{eq:J-small}
K^{-2}\E_P[\1_{E_{K,M}}J_K]\longrightarrow0.
\end{equation}
\end{lemma}

\begin{proof}
For deterministic ordered $z_1,\ldots,z_K\in(0,\ell)$, let $A(t)$ be
their counting function and $\rho=K/\ell$. Direct expansion yields the
exact identity
\begin{equation}\label{eq:quantile}
\int_0^\ell(A(t)-\rho t)^2\dd t
=\rho\sum_{i=1}^K
\left(z_i-\frac{i-\frac12}{\rho}\right)^2+\frac{\ell}{12}.
\end{equation}
For example,
$\int_0^\ell A(t)^2\dd t=\sum_i(2i-1)(\ell-z_i)$ and
$\int_0^\ell tA(t)\dd t=\frac12\sum_i(\ell^2-z_i^2)$, which verify
\eqref{eq:quantile}.

Take $z_i=x_i-a$ and put $D_0(s)=N_{(0,s]}-s$. Except at endpoints of
Lebesgue measure zero,
\[
A(t)-\rho t=D_0(a+t)+(a-1)+(1-\rho)t.
\]
On $E_{K,M}$, $a\le M$, $\ell\asymp K$, and $b\le3K$ for all
sufficiently large $K$ (with $M$ fixed). The midpoint locations
$(i-\frac12)\ell/K$ differ from $i\ell/L$ by at most a fixed constant.
It follows from \eqref{eq:quantile} that, on this event,
\begin{equation}\label{eq:quadratic-bound}
\sum_i(u_i-i)^2
\le C\int_0^{3K}D_0(s)^2\dd s+C_MK+CK(\ell-L)^2.
\end{equation}
Now $\int_0^{3K}v_P(s)\dd s=o(K^2)$ by \eqref{eq:hyper}, and the last
term on the good event is at most $Cr_K^2K^2$. Take expectations.
Because the good event is retained-exterior measurable, conditional
expectation turns the left side into the expression in
\eqref{eq:J-small}.
\end{proof}

\section{Tightness of the exterior force cost}\label{sec:forces}

The force estimates will only use tightness. We will not require an
integrable product of an exterior energy and an interior count envelope.

\Needspace{13\baselineskip}
\begin{lemma}[A half-line force envelope]\label{lem:halfline}
At a particle endpoint, remove the endpoint itself and count residual
exterior points at outward distance at most $t$. For a density-one
configuration satisfying Lemma~\ref{lem:as-regularity}, put
\[
D(t)=N_{\rm ext}(t)-(t-1)_+,\qquad
H(x)=\int_0^\infty\frac{|D(t)|}{(x+t)^2}\dd t,\quad x\ge0.
\]
Then
\begin{equation}\label{eq:B}
\calB=H(0)^2+\int_0^\infty H(x)^2\dd x<\infty.
\end{equation}
After an affine scale $c\in[1/2,2]$, the analogous quantity centered at
the new density $\lambda=1/c$ is bounded by $C(\calB+1)$.
\end{lemma}

\begin{proof}
Simplicity and local finiteness give a positive distance to the next
residual exterior point. Thus $D(t)=0$ for all sufficiently small $t$.
At infinity,
$\int D(t)^2t^{-2}\dd t<\infty$ by
Lemma~\ref{lem:as-regularity}. Cauchy--Schwarz proves $H(0)<\infty$.
Applying the Schur test to the kernel $t/(x+t)^2$, with weight
$t^{-1/2}$, gives
\begin{equation}\label{eq:hardy}
\int_0^\infty H(x)^2\dd x
\le\frac{\pi^2}{4}\int_0^\infty D(t)^2t^{-2}\dd t.
\end{equation}

For the scaled configuration the centered discrepancy is
\[
D(t/c)+(t/c-1)_+-c^{-1}(t-1)_+.
\]
The last two terms differ by a uniformly bounded function vanishing
for $t<1/2$. Its force envelope is at most $C/(x+1/2)$.
A change of variables in the first term proves the scaling assertion.
\end{proof}

Let $B_\zeta(u)$ be the difference of the actual and lattice exterior
forces in the rescaled block. More precisely, with endpoint terms
cancelled,
\[
B_\zeta(u)=\PV\sum_{p\in c(\zeta-a)\setminus\{0,L\}}\frac1{u-p}
-\PV\sum_{p\in\xi_K^{\rm lat}\setminus\{0,L\}}\frac1{u-p}.
\]
For almost every retained exterior under a candidate law, each paired
sum exists and their difference is bounded on $[0,L]$. Define
\begin{equation}\label{eq:Dcost}
\mathcal D_K(\zeta)=
\E_{\mu_\zeta}\sum_{i=1}^K B_\zeta(u_i)^2.
\end{equation}

\begin{proposition}[Tight conditional force cost]\label{prop:force}
For fixed $M$, the nonnegative random variables
$\1_{E_{K,M}}\mathcal D_K$ are tight as $K\to\infty$.
\end{proposition}

\begin{proof}
On the good event, $c=L/\ell\in[1/2,2]$ for large $K$. The scaled
exterior density is $\lambda=\ell/L$. Center each scaled outward
counting function at $\lambda(t-1)_+$.
Stieltjes integration by parts identifies its force difference from
the corresponding continuum ray with an integral of $D(t)/(u+t)^2$.
The boundary terms vanish by density one.

For the unit lattice the residual outward count is $\lfloor t\rfloor$;
its difference from $(t-1)_+$ is bounded and vanishes below $1$.
The continuum density mismatch between the two paired rays integrates
explicitly to a signed version of
$(\lambda-1)\log((u+1)/(L-u+1))$. Consequently,
\begin{equation}\label{eq:force-pointwise}
\begin{split}
|B_\zeta(u)|\le{}&
H_+(u)+H_-(L-u)+\frac C{u+1}+\frac C{L-u+1}\\
&+|\lambda-1|\abs{\log\frac{u+1}{L-u+1}}.
\end{split}
\end{equation}
The endpoint charges at $0,L$ cancel exactly before this estimate.
Here the half-line quantities $\calB_\pm$ depend only on the retained
exterior.

For any decreasing nonnegative $h$ and interior configuration,
integration by parts gives
\begin{equation}\label{eq:sum-envelope}
\sum_i h(u_i)^2
\le C_{\rm int}^+
\left(h(0)^2+\int_0^L h(s)^2\dd s\right),\qquad
C_{\rm int}^+=\sup_{t\ge0}\frac{\#\{i:u_i\le t\}}{1+t}.
\end{equation}
For instance use $N_{\rm int}(t)\le C_{\rm int}^+(1+t)$ in
$\int h^2\dd N_{\rm int}$; the terminal boundary term gives exactly
the expression on the right. The right endpoint has the reflected
bound with $C_{\rm int}^-$.

The function $h(s)=\log((L+1)/(s+1))$ satisfies
$h(0)^2+\int_0^Lh(s)^2\dd s\le CL$.
Squaring \eqref{eq:force-pointwise}, using
\eqref{eq:sum-envelope}, and taking conditional expectations gives
\begin{equation}\label{eq:force-total}
\mathcal D_K\le
C\{X_+(\calB_++1)+X_-(\calB_-+1)\}
+C(\lambda-1)^2L(X_++X_-),
\end{equation}
where $X_\pm=\E_{\mu_\zeta}C_{\rm int}^\pm$.

It remains to verify tightness under the random boundary law.
Lemma~\ref{lem:halfline} gives $\calB_\pm<\infty$ Palm-almost surely.
Equation~\eqref{eq:first-positive}, restricted to $\gap_*\le M$, and
Palm point-shift invariance bound the laws of the endpoint environments
by $M P^0$, uniformly in $K$. The scaling bound in
Lemma~\ref{lem:halfline} therefore makes $\calB_\pm$ uniformly tight on
this event.

The rescaled interior count envelopes are at most twice the corresponding
full half-line envelopes of the unscaled configuration. Since
$E_{K,M}$ is exterior measurable, conditional expectation and
Lemma~\ref{lem:palm-envelope} give
\[
\E_P[\1_{E_{K,M}}X_\pm]\le
2M\E_{P^0}C^\pm<\infty.
\]
Thus $\1_{E_{K,M}}X_\pm$ are tight as well. Products of tight
nonnegative variables are tight, without any independence assumption.
Finally,
\[
(\lambda-1)^2L=\frac{(\ell-L)^2}{L}\le Cr_K^2
\]
on $E_{K,M}$. Insert these facts into \eqref{eq:force-total}.
\end{proof}

\section{Proof of the discrepancy comparison theorem}
\label{sec:proof-main}

Let $P,Q$ satisfy Theorem~\ref{thm:comparison}. Fix $M,r,G$ and choose
$r_K$ in Lemma~\ref{lem:length} common to the two processes.
For almost every retained exterior, the logarithmic derivative of the density
ratio $\mu_\zeta/\omega_{K,T}$ has no contribution from the interior
pair interactions or from the common endpoint charges. Therefore
\begin{equation}\label{eq:score}
I(\mu_\zeta\mid\omega_{K,T})
\le C_\beta\mathcal D_K(\zeta)+\frac{2J_K(\zeta)}{T^2}.
\end{equation}
For a fixed exterior, the residual exterior forces are bounded on the
closed interval, since the next exterior particles have positive
distance from the endpoints. Relative entropy and Fisher information
are consequently finite. This also avoids an inverse-gap moment
assumption for small $\beta$.

Define the deterministic sequence
\[
q_K=\frac1{K^2}\left\{
\E_P[\1_{E_{K,M}}J_K]+\E_Q[\1_{E_{K,M}}J_K]\right\}\longrightarrow0
\]
and choose
\begin{equation}\label{eq:choice}
\varepsilon_K=\max\{\sqrt{q_K},K^{-1/2}\},
\qquad T_K=\varepsilon_KK.
\end{equation}
The same reference $\omega_{K,T_K}$ is used for both processes.
Proposition~\ref{prop:gap-entropy} and \eqref{eq:score} imply
\begin{equation}\label{eq:comparison-final}
|\mu_\zeta(F_K)-\omega_{K,T_K}(F_K)|^2
\le C_G\left\{\varepsilon_K\mathcal D_K
                 +\frac{J_K}{\varepsilon_KK^2}\right\}.
\end{equation}
On the good event, the first term tends to zero in probability by
Proposition~\ref{prop:force}; the second tends to zero in mean by
Lemma~\ref{lem:quadratic} and \eqref{eq:choice}. The left side is
bounded by $4\norm G_\infty^2$, so the expectation of its square root
over good boundaries tends to zero.

Let $\mathfrak F_K(\gamma)$ denote the rescaled block statistic
$F_K(u(\gamma))$. Conditional DLR and exterior measurability of the
good event imply
\[
\limsup_{K\to\infty}
\abs{\E_P\mathfrak F_K-\omega_{K,T_K}(F_K)}
\le 2\norm G_\infty P(\gap_*>M).
\]
The same estimate holds for $Q$. Cancelling the common reference and
then sending $M\to\infty$ proves
\begin{equation}\label{eq:cancel}
\lim_{K\to\infty}
\abs{\E_P\mathfrak F_K-\E_Q\mathfrak F_K}=0.
\end{equation}

The affine rescaling does not alter the limiting gap statistic.
Extending $G$ by zero gives a bounded globally Lipschitz function, and
each gap occurs in at most $r$ windows. Hence, with
$\mathfrak F_K^{\rm un}$ the unscaled statistic,
\[
|\mathfrak F_K-\mathfrak F_K^{\rm un}|
\le \frac{C_G}{K-r}|c-1|\ell
=\frac{C_G}{K-r}|\ell-L|\longrightarrow0
\]
almost surely by density one. Boundedness gives convergence in mean.
Lemma~\ref{lem:identify} and \eqref{eq:cancel} now yield
\[
\E_{P^0}G(g_0,\ldots,g_{r-1})
=\E_{Q^0}G(g_0,\ldots,g_{r-1})
\]
for every $r$ and $G\in C_c^2((0,\infty)^r)$.
Thus all finite consecutive-gap distributions agree. Point-shift
invariance extends this to every finite window of integer gap indices,
so the Palm laws coincide. Palm inversion at intensity one,
\[
\E_P f(\gamma)=
\E_{P^0}\left[\int_0^{g_0}f(\theta_t\gamma)\dd t\right],
\]
then gives $P=Q$. This proves Theorem~\ref{thm:comparison}.

\section{From logarithmic electric energy to discrepancy control}
\label{sec:energy}

This section supplies the finite-energy implication separately from
the comparison proof. We use the standard logarithmic regularization.
Let $\rho_\eta$ be uniform probability measure on the circle of radius
$\eta$ in $\R^2$. If
$-\diver E=2\pi(\gamma-dx\,\delta_0)$, set
\[
f_\eta(z)=\1_{\{|z|\le\eta\}}\log(|z|/\eta),\qquad
F=E_\eta=E+\sum_{p\in\gamma}\nabla f_\eta(\,\cdot-(p,0)).
\]
Then
\begin{equation}\label{eq:regularized-div}
-\diver F=2\pi(\gamma*\rho_\eta-dx\,\delta_0).
\end{equation}

\begin{remark}[Convention in the supplied version]\label{rem:cutoff}
The displayed cutoff in \cite[Definition~2.20]{DHLM} is printed with
$|z|\ge\eta$. That expression does not cancel the point singularities.
Here the cutoff is explicitly taken with $|z|\le\eta$, as required by
the regularized electric-field interpretation and by
\eqref{eq:regularized-div}. The argument below uses only finiteness of
the spatial mean energy at one fixed $\eta$, so numerical normalizations
of the renormalized-energy functional do not enter its conclusion.
\end{remark}

\begin{lemma}[Stationary electric lift]\label{lem:lift}
Let $P$ be ergodic and suppose $P$-almost every configuration has finite
renormalized electric energy. There are $\eta>0$ and a stationary
joint law of $(\gamma,F)$, with configuration marginal $P$, satisfying
\eqref{eq:regularized-div} and
\begin{equation}\label{eq:field-energy}
\calE=\E\int_{[0,1]\times\R}|F(x,y)|^2\dd x\dd y<\infty.
\end{equation}
\end{lemma}

\begin{proof}
Choose a $P$-generic configuration $\gamma$ with finite energy. No
attainment of the defining infimum over fields is needed: choose a
compatible field with finite renormalized functional. For at least
one sufficiently small fixed $\eta$, this gives
\[
A=\limsup_{R\to\infty}\frac1{2R}
\int_{[-R,R]\times\R}|E_\eta|^2<\infty.
\]
Write $F=E_\eta$ and average translates of the single pair:
\[
\mathsf Q_R=\frac1{2R}\int_{-R}^R
\delta_{(\theta_x\gamma,\theta_xF)}\dd x.
\]
Its configuration marginal converges to $P$. For each $a,b>0$,
\[
\E_{\mathsf Q_R}
\int_{[-a,a]\times[-b,b]}|F|^2
\le
\frac{2a}{2R}
\int_{[-R-a,R+a]\times\R}|F|^2.
\]
Local $L^2$ bounds give tightness in a local distribution topology,
for instance local negative Sobolev spaces where bounded $L^2$ sets
are relatively compact. Together with configuration tightness, a
diagonal subsequence has a joint limit. Translation invariance follows
because translating an averaging interval by a fixed amount changes
its normalized mass by $O(R^{-1})$. The divergence identity passes in
distributions: smearing the configuration by the fixed compactly
supported circle measure is continuous in this topology.
Lower semicontinuity on bounded vertical strips, followed by monotone
convergence, proves \eqref{eq:field-energy}. Thus no measurable
configuration-by-configuration selection of fields is required.
\end{proof}

\begin{proposition}[Spectral energy and hyperuniformity]
\label{prop:spectral}
Suppose $P$ is stationary of intensity one, satisfies
\eqref{eq:linear}, and admits a stationary electric lift as in
Lemma~\ref{lem:lift}. Then \eqref{eq:hyper} and
\eqref{eq:integrated} hold.
\end{proposition}

\begin{proof}
Let $\sigma=\gamma-dx$. Finite local second moments define its
covariance spectral measure $S$, normalized by
\[
\E_P|\langle\sigma,f\rangle|^2
=\int_\R|\widehat f(k)|^2S(\dd k),\qquad
\widehat f(k)=\int_\R e^{-ikx}f(x)\dd x.
\]
The spectral measure can be singular. For interval counts,
\begin{equation}\label{eq:count-spectrum}
v_P(t)=\int_\R\frac{4\sin^2(kt/2)}{k^2}S(\dd k),
\end{equation}
with the continuous value $t^2$ at zero.
The linear bound implies $S(\{0\})=0$.

Choose a real even smooth $\chi$ with compact spectral support away
from zero. Let $\varphi(x,y)$ be the stationary random potential
obtained by convolving $\sigma$ horizontally with multiplier
$\chi(k)e^{-|k||y|}$. Its mean is zero. Direct spectral calculation
gives the averaged gradient energy
\begin{equation}\label{eq:test-energy}
\E\int_{[0,1]\times\R}|\nabla\varphi|^2
=2\int_\R |k|\chi(k)^2S(\dd k).
\end{equation}
Pair \eqref{eq:regularized-div} with this potential, average per unit
horizontal length, and take expectation. The result is
\begin{equation}\label{eq:test-pairing}
\E\int_{[0,1]\times\R}F\cdot\nabla\varphi
=2\pi\int_\R\chi(k)a_\eta(k)S(\dd k),
\end{equation}
where
\[
a_\eta(k)=\frac1{2\pi}\int_0^{2\pi}
e^{ik\eta\cos\theta-|k|\eta|\sin\theta|}\dd\theta.
\]
For completeness, \eqref{eq:test-pairing} is a stationary
integration-by-parts identity, not a pointwise evaluation of an
$L^2$ field. First multiply the test by smooth horizontal and vertical
cutoffs and smooth its cusp at $y=0$. The distribution identity then
holds pathwise, even though the test is a function of $\gamma$.
After division by horizontal length, the horizontal cutoff terms
vanish by stationarity and Cauchy--Schwarz. The spectral support
bounded away from zero gives exponential vertical decay and finite
second moments for the test and its gradient. It therefore permits
removal of the vertical cutoff and the smoothing. The deterministic
background term, including the deterministic difference between the
smeared line and the original line, has zero expected pairing because
the test has mean zero. The remaining covariance pairing is the
multiplier on the right of \eqref{eq:test-pairing}.

Equations \eqref{eq:field-energy}--\eqref{eq:test-pairing} and
Cauchy--Schwarz imply, with a numerical constant $C$,
\begin{equation}\label{eq:spectral-duality}
\abs{\int\chi a_\eta\dd S}^2
\le C\calE\int |k|\chi^2\dd S.
\end{equation}
The function $a_\eta$ is real, continuous and equals one at zero.
Choose $\delta>0$ such that $a_\eta\ge1/2$ on $|k|<\delta$.
For a smooth even cutoff $0\le\psi\le1$ supported in
$\varepsilon<|k|<\delta$, take
$\chi=a_\eta\psi/|k|$. Smooth approximation within this annulus is
equivalent if needed. Writing
\[
A_\psi=\int\frac{a_\eta(k)^2}{|k|}\psi(k)S(\dd k),
\]
the right-hand integral in \eqref{eq:spectral-duality} is at most
$A_\psi$, since $\psi^2\le\psi$. Thus $A_\psi\le C\calE$.
Exhausting the punctured interval proves
\begin{equation}\label{eq:low-frequency}
\int_{0<|k|<\delta}\frac{S(\dd k)}{|k|}<\infty.
\end{equation}
This cutoff argument does not assume the desired integrability.

The high-frequency estimate follows from
$\int_1^2v_P(t)\dd t<\infty$ and \eqref{eq:count-spectrum}.
For fixed $\delta>0$, the continuous function
$\int_1^2\sin^2(kt/2)\dd t$ has a strictly positive lower bound on
$|k|\ge\delta$: it is positive on every compact annulus and tends
to $1/2$ as $|k|\to\infty$. Hence
\begin{equation}\label{eq:high-frequency}
\int_{|k|\ge\delta}k^{-2}S(\dd k)<\infty.
\end{equation}

Tonelli in \eqref{eq:count-spectrum}, followed by the elementary bound
\[
\frac4{k^2}\int_1^\infty\frac{\sin^2(kt/2)}{t^2}\dd t
\le
\begin{cases}
C/|k|,&0<|k|<\delta,\\
4/k^2,&|k|\ge\delta,
\end{cases}
\]
proves \eqref{eq:integrated}. Finally, the integrand in
$v_P(t)/t$ tends to zero for every $k\ne0$ and has the same two
integrable majorants for $t\ge1$. Dominated convergence proves
\eqref{eq:hyper}.
\end{proof}

\begin{proof}[Proof of Theorem~\ref{thm:finite-energy}]
Lemma~2.23 of \cite{DHLM} gives
\[
\E_P(N_I-|I|)^2\le C(C+\E_PW)|I|.
\]
Stationarity then forces intensity one: otherwise the square of the
mean discrepancy would grow quadratically in $|I|$. The same estimate
gives \eqref{eq:linear}.
Proposition~3.13 of \cite{DHLM} decomposes $P$ into ergodic stationary
canonical DLR solutions with finite expected energy. For each component,
Lemma~\ref{lem:lift} and Proposition~\ref{prop:spectral} give the remaining
conditions of Theorem~\ref{thm:comparison}. All such components therefore
coincide.

The existence theorem of \cite{DHLM} supplies the stationary finite-energy
solution $\Sine_\beta$. Apply the same decomposition to it. Its
components belong to the same singleton class, so their common law is
$\Sine_\beta$. Every candidate $P$, being a mixture of that law, equals
$\Sine_\beta$.
\end{proof}

\section{Relation to conditional universality}

At $\beta=2$, the deterministic theorem of \cite[Theorem~1.3]{KMD}
provides another route. The reciprocal principal value and density-one
conditions needed there follow from the discrepancy bounds, so its
kernel convergence applies to every admissible candidate boundary.
DLR and convergence of compactly supported Laplace functionals then
identify the process.

For general $\beta$, the main theorem of the supplied version of
\cite{BEY} concerns a fixed real analytic uniformly convex potential
on the line, and its local comparison is combined with strong
particle-location estimates. A direct application to arbitrary
random DLR exterior potentials would require checking assumptions not
provided by macroscopic potential convergence alone.
The present proof uses the static convexity inequalities from that
framework and compares with one common reference law. In particular,
it establishes neither a separate universality theorem for all
deterministic boundaries nor a limiting theorem for
$\omega_{K,T_K}$.

The distinction between the two energy inputs is useful:
$v_P(t)=o(t)$ makes the average quadratic displacement $o(K)$ per
particle, whereas $\int_1^\infty v_P(t)t^{-2}\dd t<\infty$ supplies
the half-line force envelopes. The linear bound alone gives neither
of these two improvements. The weak confinement can be chosen between
the quadratic displacement scale and the block length precisely
because the improvement is little-$o$; no fixed power rate is needed.

\begin{thebibliography}{DHLM19}

\bibitem[BEY12]{BEY}
Paul Bourgade, L\'aszl\'o Erd\H{o}s and Horng-Tzer Yau,
\emph{Universality of General $\beta$-Ensembles}.
Supplied preprint arXiv:1104.2272v5, 5 February 2012.
All theorem and equation references in this manuscript refer to that
version. The bibliography of \cite{DHLM} records the published article
in \emph{Duke Mathematical Journal} \textbf{163} (2014), no.~6,
1127--1190.

\bibitem[DHLM19]{DHLM}
David Dereudre, Adrien Hardy, Thomas Lebl\'e and Myl\`ene Ma\"ida,
\emph{DLR equations and rigidity for the Sine-beta process}.
Supplied preprint arXiv:1809.03989v2, arXiv version date
15 November 2019; title-page date 18 November 2019.

\bibitem[KMD19]{KMD}
Arno B.~J. Kuijlaars and Erwin Mi\~na-D\'iaz,
\emph{Universality for conditional measures of the sine point process}.
\emph{Journal of Approximation Theory} \textbf{243} (2019), 1--24.
Supplied preprint arXiv:1703.02349v2, 21 March 2019.

\end{thebibliography}
\end{document}
